\documentstyle[%  


righttag%  


,11pt%  


,amssymb%  


,russian%               remove in Eng.  


,izvru%                 Set izven% in Eng.  


]{amsart}  


  


% {=}


 


\newcommand{\la}{\langle} 


\newcommand{\ra}{\rangle} 


\newcommand{\ch}{\operatorname{ch}}  


\newcommand{\sh}{\operatorname{sh}}  


\newcommand{\re}{\operatorname{Re}}  


\newcommand{\loc}{\operatorname{loc}}  


\newcommand{\th}{\operatorname{th}}  


\newcommand{\tts}[1]{}  


  


\theoremstyle{plain} 


\theorembodyfont{\it}


\newtheorem{thmnonum}{\hskip\parindent Теорема}


\renewcommand{\thethmnonum}{}





%\hoffset=-15mm%                  For Eng.  


\setcounter{page}{65}


\god{2001}


\nomer{4~(467)} %                        Remove in Eng.  


%\nomer{Vol.\,45, No.\,4}%               Set in Eng.  


%\str{\pageref{begin}--\pageref{end}}%  Set in Eng.  


\udk{517.518}  


  


\author{\it Вит.\,В.\,Волчков}  


% NB:   ^^ remove in Eng.  


\title{Локальная теорема о двух радиусах\\ для $\cal M$-гармонических 


функций} 


  


\date{}  


% \pagestyle{myheadings}%         Set in Eng.  


\pagestyle{plain} %              Remove in Eng.  


\begin{document}  


% \thispagestyle{plain}%          Set in Eng.  


\maketitle  


\label{begin}  


 


\section*{Введение} 


 


Пусть $f$ --- непрерывная функция на вещественном евклидовом 


пространстве 


$R^n$, $n\ge2$, $E$ --- заданное множество положительных чисел. 


Предположим, что при всех $r\in E$ и $x\in R^n$ 


\begin{equation} 


f(x)=\frac{1}{\omega r^n}\int_{|y|\le r}f(x+y)d y, 


\end{equation} 


где $\omega$ --- объем единичного шара в $R^n$ и $|\cdot|$ --- 


евклидова норма 


в $R^n$. Для каких $E$ отсюда следует, что $f$ --- гармоническая 


функция? Известная теорема Дельсарта о двух радиусах утверждает, что 


функция 


$f$ гармоническая, если $E$ состоит из двух чисел $r_1$ и 


$r_2$, отношение которых не является отношением корней целой функции 


$2^{n/2}\Gamma\left(\frac{n}{2}+1\right)\Dot J_{n/2}(z)/z^{n/2}-1$ (см. 


[1], а также обзоры [2]--[4] с обширной библиографией). 


Примеры показывают (напр., [5]), что указанное 


условие для $r_1/r_2$ является необходимым. 


 


Получение локального аналога теоремы Дельсарта (напр., в случае, когда 


функция с условием (1) задана в шаре $\cal B_R=\{x\in R_n:|x|<R\}$ 


радиуса $R>\max(r_1,r_2)$) является существенно более трудной задачей. 


Здесь возникают новые эффекты: наличие гармоничности во многом зависит 


не 


только от природы чисел $r_1$ и $r_2$, но и от размеров $\cal B_R$. 


Локальная 


теорема о двух радиусах при условии $R>r_1+r_2$ получена различными 


методами в [6] и [7]. В [7] 


доказана также точность этого условия на размеры $\cal B_R$. Более 


того, 


дальнейшее развитие техники, предложенной в [7], позволило полностью 


изучить вопрос 


о гармоничности функции $f$ с условием (1) и при всех $R\le r_1+r_2$


[5]. 


Большой интерес представляет обобщение этих результатов на некомпактные 


двухточечно-однородные пространства (в соответствии с их классификацией 


(напр., [8], с.\,203) --- это в точности евклидовы пространства, 


гиперболические пространства $H^n(R)$, $H^n(\Bbb C)$, $H^n(\Bbb H)$ и 


гиперболическое 


пространство Кэли $H^{16}(Cay)$). Первые результаты в этом направлении 


принадлежат К.А.\,Беренстейну и Л.\,Зальцману [9], 


которые распространили теорему Дельсарта на указанный 


класс пространств. За исключением евклидова случая локальный 


аналог этой теоремы был получен лишь для вещественного гиперболического 


пространства $H^n(R)$ ([10], теорема 8). В данной работе доказана 


локальная теорема о двух радиусах для гармонических функций на 


комплексном 


гиперболическом пространстве $H^n(\Bbb C)$. 


 


\section{Формулировка основного результата} 


 


Пусть $\Bbb C^n$ --- комплексное евклидово пространство размерности 


$n\ge2$ 


с эрмитовым скалярным произведением $\la\cdot,\cdot\ra$, $B=\{z\in\Bbb 


C^n:|z|<1\}$, 


где $|z|^2=\la z,z\ra$. Обозначим через $d(z,w)$ расстояние между 


точками 


$z,w\in B$ в метрике Бергмана, т.\,е. 


\begin{equation} 


d(z,w)=\frac{1}{2}\ln\bigg( 


\frac{|1-\la w,z\ra|+\sqrt{|w-z|^2+ 


|\la w,z\ra|^2-|w|^2|z|^2}} 


{|1-\la w,z\ra|-\sqrt{|w-z|^2+|\la w,z\ra|^2- 


|w|^2|z|^2}}\bigg).


\end{equation} 


Как известно (напр., [11]), комплексное гиперболическое пространство 


$H^n(\Bbb C)$ размерности $n$ изометрично шару $B$ с метрикой (2). 


Риманова мера на $H^n(\Bbb C)$ имеет вид 


$d\mu(z)=d m(z)/(1-|z|^2)^{n+1}$, 


где $d m$ --- мера Лебега в $\Bbb C^n$. Оператор Лапласа--Бельтрами на 


$H^n(\Bbb C)$ 


имеет вид ([11]; [12], с.\,54) 


$$ 


\wt\Delta=4(1-|z|^2)\sum^n_{i,j=1} 


(\delta_{i j}-z_i \ov z_j) 


\frac{\partial^2}{\partial z_i \partial\ov z_j}, 


$$ 


где $\delta_{i j}$ --- символ Кронекера. Функция $f\in C^2(B)$ 


называется 


$\cal M$-гармонической в $B$, если $\wt\Delta f=0$. 


 


Пусть $\alpha,\beta\ge0$, $\lambda\in\Bbb C$, 


$\varphi^{(\alpha,\beta)}_\lambda(t)=F\left( 


\frac{\alpha+\beta+1-i\lambda}{2},\ 


\frac{\alpha+\beta+1+i\lambda}{2};\ \alpha+1;\ -(\sh t)^2\right)$, 


где $F$ --- гипергеометрическая функция. При фиксированном $r>0$ 


обозначим через $M(r)$ множество корней~$\lambda$ функции 


$\eta_r(\lambda)=1-\ch^2r\cdot\varphi^{(n,1)}_\lambda(r)$, 


лежащих в полуплоскости $\re\lambda\ge0$ и не принадлежащих 


отрицательной части мнимой оси $\{it,\ t<0\}$. Пусть также 


$M(r_1,r_2)=M(r_1)\cap M(r_2)$, $B_R=\{z\in B:d(z,0)<R\}$. 


 


Основным результатом данной работы является 


 


\begin{thmnonum} 


Пусть $r_1,r_2>0$, $R>\max(r_1,r_2)$, $f\in C(B_R)$ и 


\begin{equation} 


f(z)=\frac{1}{\mu(B_{r_j})}\int_{d(z,w)\le r_j} 


f(w)d\mu(w) 


\end{equation} 


при всех $z\in B_{R-r_j}$, $j=1,2$. Тогда 


 


\begin{itemize} 


\item[1)] eсли $r_1+r_2<R$ и $M(r_1,r_2)=\{in\}$, то $f$\; $\cal 


M$-гармонична в $B_R;$


\item[2)] eсли $r_1+r_2>R$ или $M(r_1,r_2)\ne\{in\}$, то существует 


$f\in C^\infty(B_R)$ с условием $(3)$, которая не является $\cal 


M$-гармонической в $B_R$. 


\end{itemize} 


\end{thmnonum} 


 


Отметим, что если в условии теоремы предполагать $f\in 


C^\infty(B_R)$, 


то ее первое утверждение верно и при $r_1+r_2\le R$, 


$M(r_1,r_2)=\{in\}$. 


Другие результаты по аналогичным проблемам средних значений содержатся 


в [4]. 


 


\section{Обозначения и вспомогательные утверждения} 


 


Следуя ([8], с.\,326), обозначим через $T_1\times T_2$ свертку двух 


распределений 


на $H^n(\Bbb C)$, одно из которых имеет компактный носитель. Положим 


$\psi_r=\delta-\chi_r/\mu(B_r)$, где $\delta$ --- дельта-распределение 


в нуле 


пространства $H^n(\Bbb C)$, $\chi_r$ --- характеристическая функция 


(индикатор) 


шара $B_r$. Класс функций $f\in C^\infty(B_R)$, для которых 


$(f\times\psi_r)(z)=0$, $z\in B_{R-r}$, 


обозначим $\cal H_r(B_R)$. 


 


Пусть $S=\{z\in\Bbb C^n:|z|=1\}$, $\rho$, $\sigma$ --- полярные 


координаты в 


$\Bbb C^n$ ($\forall z\in\Bbb C^n$\; $\rho=|z|$, а если $z\ne0$, то 


$\sigma=z/|z|$), 


$\{S^{(k)}_{p,q}(\sigma)\}$, $1\le k\le Q(n,p,q)$, --- фиксированный 


ортонормированный базис 


в пространстве сферических гармоник бистепени $(p,q)$ ([12], с.\,265). 


Всякой функции $f\in L_{\loc}(B_R)$ соответствует ряд Фурье 


$$ 


f(z)\sim\sum^\infty_{p,q=0}\, 


\sum^{Q(n,p,q)}_{k=1}f_{p,q,k}(\rho)S^{(k)}_{p,q}(\sigma), 


\quad \rho\in(0,\th R), 


$$ 


где 


$f_{p,q,k}(\rho)=\int\limits_{S}f(\rho,\sigma)


\ov{S^{(k)}_{p,q}(\sigma)}d\sigma$. 


 


Леммы 1--6, приводимые ниже, доказываются методами работы [13]. 


 


\begin{lem} 


Пусть $f\in\cal H_r(B_R)$. Тогда функция 


$f_{p,q,k}(\rho)S^{(l)}_{p,q}(\sigma)$ принадлежит $\cal H_r(B_R)$ при 


всех $p,q\ge0$, $1\le k,l\le Q(n,p,q)$. 


\end{lem} 


 


\begin{lem} 


Пусть $f\in\cal H_r(B_R)$ и $f=0$ в $B_r$. Тогда $f=0$ в $B_R$. 


\end{lem} 


 


Пусть $\lambda\in\Bbb C$, 


$\gamma=\gamma(\lambda)=\frac{n-i\lambda}{2}$. При $\rho\in[0,1)$, 


$m=0,1,\dotsc$ 


обозначим 


\begin{gather*} 


\Phi^{p,q}_\lambda(\rho)= 


\frac{\Gamma(\gamma+q)\Gamma(\gamma+p)\Gamma(n)} 


{\Gamma(n+p+q)\Gamma^2(\gamma)}\rho^{p+q} 


(1-\rho^2)^\gamma F(\gamma+q,\gamma+p;n+p+q;\rho^2),\\ 


% 


\Phi^{p,q}_{\lambda,m}(\rho)=\frac{d^m}{d\mu^m} 


\big(\Phi^{p,q}_\mu(\rho)\big)\bigg|_{\mu=\lambda}, 


\end{gather*} 


где $\Gamma$ --- гамма-функция. Сферические функции на $H^n(\Bbb C)$ 


имеют вид [12] 


$\varphi_\lambda(z){=}\varphi^{(n-1,0)}_\lambda(d(0,z)){=}\linebreak


\Phi^{0,0}_\lambda(|z|)$. Положим 


$\varphi_{\lambda,m}(z)=\Phi^{0,0}_{\lambda,m}(|z|)$. 


 


\begin{lem} 


Для любого $z\in B$ 


\begin{equation} 


(\varphi_{\lambda,m}\times\psi_r)(z)=\sum^m_{j=0} 


\frac{m!}{j!(m-j)!}\eta^{(j)}_r(\lambda) 


\varphi_{\lambda,m-j}(z).


\end{equation} 


\end{lem} 


 


\begin{lem} 


Пусть $n_\lambda$ --- кратность корня $\lambda\in M(r)$, 


$m=0,\dotsc,n_\lambda-1$. 


Тогда функция $\Phi^{p,q}_{\lambda,m}(\rho)S^{(k)}_{p,q}(\sigma)$ 


принадлежит $\cal H_r(B_R)$ 


при всех $p,q\ge0$, $1\le k\le Q(n,p,q)$. 


\end{lem} 


 


\begin{lem} 


Пусть $f\in C^\infty(B_R)$. Тогда для того чтобы $f\in\cal H_r(B_R)$, 


необходимо и достаточно, чтобы при всех целых $p,q\ge0$ и 


$1\le k\le Q(n,p,q)$ 


имело место равенство 


$$ 


f_{p,q,k}(\rho)=\sum_{\lambda\in M(r)}\, 


\sum^{n_\lambda-1}_{m=0}c_{\lambda,m,p,q,k} 


\Phi^{p,q}_{\lambda,m}(\rho), 


$$ 


причем 


\begin{equation} 


\max_{0\le m\le n_\lambda-1}|c_{\lambda,m,p,q,k}|= 


O(|\lambda|^B) 


\end{equation} 


при $|\lambda|\to+\infty$ и любом фиксированном $B<0$.


\end{lem} 


 


\begin{lem} 


Пусть $r_1,r_2>0$, $r_1+r_2\le R$, $M(r_1,r_2)=\{in\}$. 


Тогда, если функция 


$$ 


f(z)=\sum_{\lambda\in M(r_1)}\, 


\sum^{n_\lambda-1}_{m=0}c_{\lambda,m} 


\Phi^{p,q}_{\lambda,m}(\rho)S^{(k)}_{p,q}(\sigma), 


\quad z\in B_R, 


$$ 


принадлежит $\cal H_{r_2}(B_R)$, причем $c_{\lambda,m}$ удовлетворяют 


$(5)$, 


то $f$ \; $\cal M$-гармонична в $B_R$. 


\end{lem} 


 


\section{Схема доказательства теоремы} 


 


1. Пусть сначала $f\in C^\infty(B_R)$. Используя 


леммы 1, 5 и 6, получаем, что все коэффициенты Фурье функции $f$\; 


$\cal M$-гармоничны в $B_R$. Это означает, что $f$ является $\cal 


M$-гармонической 


в $B_R$. Общий случай сводится к рассмотренному стандартным 


приемом сглаживания (напр., [8], доказательство теоремы 2.6, с. 125). 


 


2. В случае, когда $M(r_1,r_2)\ne\{in\}$, всем 


требованиям утверждения 2 удовлетворяет функция $\varphi_\lambda(z)$, 


где 


$\lambda\in M(r_1,r_2)\setminus\{in\}$. Пусть $r_1<r_2$, $r_1+r_2>R$ и 


$M(r_1,r_2)=\{in\}$. Зафиксируем $\varepsilon\in(0,\min\{r_1+r_2-R,\ 


R-r_2\})$. 


Рассмотрим радиальную функцию $h(\rho)$ со следующими свойствами: 


$h\not\equiv0$, $h\in C^\infty(B_{r_2+\varepsilon})$; $h(\rho)=0$ при 


$\rho\in(0,\th(R-r_1))\cup(\th(r_2-\varepsilon), 


\th(r_2+\varepsilon))$; 


$\int\limits_{R-r_1<d(z,0)<r_2-\varepsilon}h(|z|)d\mu(z)=0$. 


Тогда $h\in\cal H_{r_2}(B_{r_2+\varepsilon})$ и по лемме 5 имеем 


$$ 


h(|z|)=\sum_{\lambda\in M(r_2)}\, 


\sum^{n_\lambda-1}_{m=0}c_{\lambda,m} 


\varphi_{\lambda,m}(z), 


$$ 


где $c_{\lambda,m}$ удовлетворяют (5). Заменяя, если необходимо, 


функцию $h$ 


подходящей линейной комбинацией вида 


$\alpha_0h+\alpha_1\wt\Delta h+\dotsb+\alpha_k\wt\Delta^k h$, 


можно считать, что $c_{\lambda,m}=0$ при 


$|\lambda|\le\varkappa=\min\{t\ge n+1:n_\lambda=1,\ |\lambda|\ge t\}$. 


Это означает, что 


$h(|z|)=\sum\limits_{\lambda\in M(r_2),|\lambda|>\varkappa} 


c_\lambda\varphi_\lambda(z)$, 


причем $c_\lambda=O(|\lambda|^b)$ при $|\lambda|\to+\infty$ и любом 


$b<0$. Рассмотрим 


функцию 


$f(z)=\sum\limits_{\lambda\in M(r_2),|\lambda|>\varkappa} 


\frac{c_\lambda}{\eta_{r_1}(\lambda)}\varphi_\lambda(z)$, $z\in B_R$. 


Из условия $r_1<r_2$ и асимптотического поведения 


$\varphi^{(n,1)}_\lambda(r_1)$ при 


$|\lambda|\to+\infty$ ([11], доказательство леммы 4.9) имеем 


$\varphi^{(n,1)}_\lambda(r_1)\to0$ 


при $|\lambda|\to+\infty$, $\lambda\in M(r_2)$. Учитывая (4), отсюда и 


из леммы 5 


получаем $(f\times\psi_{r_1})(z)=h(|z|)$ и $f\in\cal H_{r_2}(B_R)$. 


Следовательно, функция $f$ удовлетворяет всем требованиям утверждения 


2. 


 


Автор благодарит профессора В.В.\,Волчкова за многочисленные полезные 


обсуждения. 


 


\begin{thebibliography}{99} 


 


\bibitem{1} 


Delsarte J. {\em Note sur une propriete nouvelle des fonctions 


hormoniques} 


// C.\,R. Acad. Sci. 


Paris Ser.A--B. -- 1958. -- V.\,246. -- \No\,9. -- P.\,1358--1360. 


\bibitem{2} 


Беренстейн К.А., Струппа Д. {\em Комплексный анализ и уравнения в 


свертках} // 


Итоги науки и техн. Современ. пробл. матем. -- 1989. -- Т.\,54. 


-- С.\,5--111. 


\bibitem{3} 


Zalcman L. {\em A bibliographic survey of the Pompeiu problem} // 


Аpproxim. 


solutions of partial different. equat. -- 1992. -- P.\,185--194. 


\bibitem{4} 


Netuka I., Vesely J. {\em Mean value property and harmonic functions} 


// Classical and Modern potential Theory and Applications. -- 1994. 


-- P.\,359--398. 


\bibitem{5} 


Волчков В.В. {\em Окончательный вариант теоремы о среднем в теории 


гармонических функций} // Матем. заметки. -- 1996. -- Т.\,59. 


-- \No\,3. -- С.\,351--358. 


\bibitem{6} 


Berenstein С.A., Gay R. {\em A local version of the two-circles 


theorem} 


// Israel J. Math. -- 1986. -- V.\,55. -- \No\,3. -- P.\,267--288. 


\bibitem{7} 


Волчков В.В. {\em Теоремы о двух радиусах на ограниченных областях 


евклидовых 


пространств} // Дифференц. уравнения. -- 1994. -- 


Т.30. -- \No\,10. -- С.\,1719--1724. 


\bibitem{8} 


Хелгасон С. {\em Группы и геометрический анализ}. -- 


М.: Мир, 1987. -- 735 с. 


\bibitem{9} 


Berenstein С.A., Zalcman L. {\em Pompeiu's problem on symmetric spaces} 


// Comment. Math. Helv. -- 1980. -- V.\,55. -- \No\,4. -- P.\,593--621. 


\bibitem{10} 


Волчков В.В. {\em Проблемы типа Помпейю на многообразиях} // Докл. АН 


Украины. -- 1993. -- \No\,11. -- С.\,9--12. 


\bibitem{11} 


Harchaoui M. {\em Inversion de la transformation de Pompeiu locale dans 


les espaces hyperboliques reel et complete} // J. Anal. Math. -- 


1995. -- V.\,67. -- P.\,1--37. 


\bibitem{12} 


Рудин У. {\em Теория функций в единичном шаре из $\Bbb C^n$}. -- 


М.: Мир, 1984. -- 455 с. 


\bibitem{13} 


Волчков Вит.\,В. {\em О функциях с нулевыми шаровыми средними на 


комплексных 


гиперболических пространствах} // Матем. заметки. -- 


2000. -- Т.\,68. -- Вып.\,4. -- С.\,504--512. 


 


\end{thebibliography} 


 


\vskip 2cm 


 


\post{\it Донецкий государственный университет }{\it Поступила} 


\post{}{07.02.2000 } 


 


\label{end} 


\end{document} 











